% Part: first-order-logic
% Chapter: syntax-and-semantics
% Section: substitution

\documentclass[../../../include/open-logic-section]{subfiles}

\begin{document}

\olfileid{fol}{syn}{sub}

\olsection{Substitutie}

\begin{defn}[Substitutie in een term]
We definiëren $\Subst{s}{t}{x}$, het resultaat van het
\emph{substitueren} van $t$ voor ieder voorkomen van~$x$ in $s$,
recursief:
\begin{enumerate}
\item \indcase{s}{c}{$\Subst{\indfrm}{t}{x}$ gewoon $s$ is.}

\item \indcase{s}{y}{$\Subst{\indfrm}{t}{x}$ eveneens gewoon~$s$ is,
  mits $y$ een variabele is en $y \not\ident x$.}

\item \indcase{s}{x}{$\Subst{\indfrm}{t}{x}$ gelijk is aan~$t$.}

\item\indcase{s}{\Atom{f}{t_1, \dots, t_n}}{$\Subst{\indfrmp}{t}{x}$ gelijk is aan
  $\Atom{f}{\Subst{t_1}{t}{x}, \dots, \Subst{t_n}{t}{x}}$.}
\end{enumerate}
\end{defn}

\begin{defn}
Een term~$t$ is \emph{!!{free for}} $x$ in $!A$ als geen van de vrije
voorkomens van~$x$ in $!A$ binnen het bereik ligt van een kwantor die
een in~$t$ voorkomende variabele bindt.
\end{defn}

\begin{ex} ~
\begin{enumerate}
\item $\Obj v_8$ is vrij voor $\Obj v_1$ in $\lexists[\Obj
  v_3]\Atom{\Obj A^2_4}{\Obj v_3,\Obj v_1}$

\item $\Obj f^2_1(\Obj v_1, \Obj v_2)$ is \emph{niet} vrij voor $\Obj
  v_0$ in $\lforall[\Obj v_2]\Atom{\Obj A^2_4}{\Obj v_0,\Obj v_2}$
\end{enumerate}
\end{ex}

\begin{defn}[Substitutie in !!a{formula}]
Als $!A$ !!a{formula} is, $x$~!!a{variable} is en $t$~een term is die
!!{free for}~$x$ in~$!A$ is, dan is $\Subst{!A}{t}{x}$ het resultaat
van het substitueren van $t$ voor alle vrije voorkomens van~$x$
in~$!A$.
\begin{enumerate}
\tagitem{prvFalse}{\indcase{!A}{\lfalse}{$\Subst{\indfrm}{t}{x}$ gelijk is aan
    $\lfalse$.}}{}

\tagitem{prvTrue}{\indcase{!A}{\ltrue}{$\Subst{\indfrm}{t}{x}$ gelijk is aan
    $\ltrue$.}}{}

\item \indcase{!A}{\Atom{P}{t_1,\dots,
    t_n}}{$\Subst{\indfrm}{t}{x}$ gelijk is aan $\Atom{P}{\Subst{t_1}{t}{x},
    \dots, \Subst{t_n}{t}{x}}$.}

\item \indcase{!A}{\eq[t_1][t_2]}{$\Subst{\indfrmp}{t}{x}$ gelijk is aan
  $\Subst{t_1}{t}{x} = \Subst{t_2}{t}{x}$.}

\tagitem{prvNot}{\indcase{!A}{\lnot !B}{$\Subst{\indfrmp}{t}{x}$ gelijk is aan
    $\lnot \Subst{!B}{t}{x}$.}}{}

\tagitem{prvAnd}{\indcase{!A}{(!B \land
    !C)}{$\Subst{\indfrmp}{t}{x}$ gelijk is aan $(\Subst{!B}{t}{x} \land
    \Subst{!C}{t}{x})$.}}{}

\tagitem{prvOr}{\indcase{!A}{(!B \lor
    !C)}{$\Subst{\indfrmp}{t}{x}$ gelijk is aan $(\Subst{!B}{t}{x} \lor
    \Subst{!C}{t}{x})$.}}{}

\tagitem{prvIf}{\indcase{!A}{(!B \lif
    !C)}{$\Subst{\indfrmp}{t}{x}$ gelijk is aan $(\Subst{!B}{t}{x} \lif
    \Subst{!C}{t}{x})$.}}{}

\tagitem{prvIff}{\indcase{!A}{(!B \liff
    !C)}{$\Subst{\indfrmp}{t}{x}$ gelijk is aan $(\Subst{!B}{t}{x} \liff
    \Subst{!C}{t}{x})$.}}{}

\tagitem{prvAll}{
\indcase{!A}{\lforall[y][!B]}{$\Subst{\indfrmp}{t}{x}$
    gelijk is aan $\lforall[y][\Subst{!B}{t}{x}]$, mits $y$ een
    andere variabele dan $x$ is; anders is $\Subst{\indfrmp}{t}{x}$
    gewoon $\indfrm$.}}{}

\tagitem{prvEx}{
\indcase{!A}{\lexists[y][!B]}{$\Subst{\indfrmp}{t}{x}$
    gelijk is aan $\lexists[y][\Subst{!B}{t}{x}]$, mits $y$ een
    andere variabele dan $x$ is; anders is $\Subst{\indfrmp}{t}{x}$
    gewoon $\indfrm$.}}{}
\end{enumerate}
\end{defn}

\begin{explain}
Merk op dat substitutie niets hoeft te veranderen: als $x$ helemaal
niet in $!A$ voorkomt, dan is $\Subst{!A}{t}{x}$ gewoon~$!A$.

De beperking dat $t$ !!{free for}~$x$ in~$!A$ moet zijn, is nodig om
gevallen als het volgende uit te sluiten. Als
$!A \ident \lexists[y][x < y]$ en $t \ident y$, dan zou
$\Subst{!A}{t}{x}$ gelijk zijn aan $\lexists[y][y < y]$. De vrije
variabele $y$ wordt bij deze substitutie door de kwantor
$\lexists[y]$ ``gevangen'', dat wil zeggen gebonden, en dat is
ongewenst. We willen bijvoorbeeld dat, als $\lforall[x][!B]$ geldt,
ook $\Subst{!B}{t}{x}$ geldt. Beschouw echter
$\lforall[x][\lexists[y][x < y]]$ (hier is $!B$ gelijk aan
$\lexists[y][x < y]$). Dit is een zin die bijvoorbeeld waar is in de
 natuurlijke getallen: voor ieder getal~$x$ bestaat er een getal~$y$
 dat groter is dan het eerste. Als we $y$ als substitutieterm voor~$x$ zouden
toelaten, kregen we
$\Subst{!B}{y}{x} \ident \lexists[y][y < y]$, en dat is onwaar. We
voorkomen dit door te eisen dat geen van de vrije variabelen in~$t$
door een kwantor in~$!A$ gebonden raakt.
\end{explain}

Om omslachtige notatie te vermijden gebruiken we vaak de volgende
conventie. Als $!A$ !!a{formula} is waarin de !!{variable}~$x$ vrij
kan voorkomen, schrijven we ook~$!A(x)$ om dit aan te geven. Als
duidelijk is welke $!A$ en~$x$ worden bedoeld en $t$ een term is
(waarvan wordt aangenomen dat hij vrij voor $x$ in $!A(x)$ is), dan
schrijven we $!A(t)$ als afkorting voor $\Subst{!A}{t}{x}$. We kunnen
bijvoorbeeld zeggen: ``$!A(t)$ heet een instantie
van~$\lforall[x][!A(x)]$.'' Daarmee bedoelen we dat, als $!A$ een
willekeurige !!{formula} is, $x$~!!a{variable} is en $t$ een term is
die vrij voor~$x$ in~$!A$ is, $\Subst{!A}{t}{x}$ een instantie van
$\lforall[x][!A]$ is.
\end{document}
